\section{古典データの量子回路への符号化}
\label{sec:data-encoding}

\subsection{データ符号化の必要性}
\label{subsec:need-for-encoding}

量子機械学習で扱う入力は，通常，実数やビット列として与えられる古典データである．
一方，量子回路が直接操作する対象は量子状態である．
したがって，古典データ$x$を量子回路へ入力するには，入力依存のユニタリ演算子
$U_{\mathrm{in}}(x)$を用いて
\begin{equation}
  \ket{\psi(x)}
  \coloneq
  U_{\mathrm{in}}(x)\ket{0}^{\otimes n}
  \label{eq:input-encoding}
\end{equation}
のように量子状態へ変換する必要がある．
この操作を\textbf{データ符号化}または\textbf{特徴量写像}という．
符号化方法によって，必要な量子ビット数，回路の深さ，表現できるデータの構造が異なる．

\subsection{基底符号化}
\label{subsec:basis-encoding}

基底符号化では，古典的な$n$ビット列$b_1b_2\cdots b_n$を計算基底状態
\begin{align*}
  b_1b_2\cdots b_n
  \longmapsto
  \ket{b_1b_2\cdots b_n}
\end{align*}
へ対応させる．
例えば，$101$は$\ket{101}$へ符号化される．
初期状態$\ket{0}^{\otimes n}$から目的の基底状態を準備するには，
ビット値が$1$である位置の量子ビットへ$X$ゲートを作用させればよい．
回路は単純であるが，$n$ビットのデータに$n$量子ビットを必要とする．

\subsection{振幅符号化}
\label{subsec:amplitude-encoding}

振幅符号化では，$2^n$次元の古典ベクトル
\begin{align*}
  \bm{x}
  \coloneq
  \begin{pmatrix}
    x_0 & x_1 & \cdots & x_{2^n-1}
  \end{pmatrix}^{\mathsf{T}}
\end{align*}
を，規格化した量子状態
\begin{equation}
  \ket{x}
  \coloneq
  \frac{1}{\|\bm{x}\|}
  \sum_{j=0}^{2^n-1}x_j\ket{j}
  \label{eq:amplitude-encoding}
\end{equation}
として表す．
少数量子ビットの振幅に高次元ベクトルを格納できる一方で，
任意の$\bm{x}$に対応する状態を準備する回路自体が複雑になる場合がある．
そのため，必要量子ビット数だけでなく，状態準備の計算量も考慮する必要がある．
